\chapter{Funções}

\section{Declaração}

Funções são declaradas com \lstinline{fn}:

\begin{lstlisting}[caption={Função básica}]
fn quadrado(x) {
  return x^2
}

display quadrado(5)    // 25
\end{lstlisting}

\subsection*{Retorno implícito}

A última expressão de uma função é retornada automaticamente:

\begin{lstlisting}
fn quadrado(x) {
  x^2    // return implícito
}
\end{lstlisting}

\subsection*{Múltiplos parâmetros}

\begin{lstlisting}
fn media(a, b) {
  (a + b) / 2.0
}

display media(10, 20)    // 15.0
\end{lstlisting}

\section{Funções como valores}

Funções são valores de primeira classe — podem ser atribuídas a variáveis, passadas
como argumento e retornadas por outras funções:

\begin{lstlisting}[caption={Função como argumento}]
fn aplicar(f, x) {
  f(x)
}

fn dobro(n) { n * 2 }

display aplicar(dobro, 7)    // 14
\end{lstlisting}

\section{Closures}

Closures são funções anônimas que capturam o ambiente onde foram definidas:

\begin{lstlisting}[caption={Closure básica}]
let multiplicador = fn(fator) {
  fn(x) { x * fator }    // captura 'fator'
}

let triplo = multiplicador(3)
display triplo(10)    // 30
\end{lstlisting}

\subsection*{Closures como argumentos}

\begin{lstlisting}
fn aplicar_dois(f, a, b) {
  f(a) + f(b)
}

let soma_quadrados = aplicar_dois(fn(x) { x^2 }, 3, 4)
display soma_quadrados    // 9 + 16 = 25
\end{lstlisting}

\section{Recursão}

Funções podem chamar a si mesmas:

\begin{lstlisting}[caption={Fatorial recursivo}]
fn fat(n) {
  if n <= 1 { return 1 }
  n * fat(n - 1)
}

display fat(10)    // 3628800
\end{lstlisting}

\begin{nota}
  \hayashi{} não otimiza recursão em cauda (\textit{tail-call optimization}).
  Para sequências longas, prefira laços.
\end{nota}

\section{Parâmetros opcionais e valores padrão}

\hayashi{} não tem parâmetros opcionais nativos, mas o padrão idiomático é usar
\lstinline{nil} como sentinela e verificar explicitamente:

\begin{lstlisting}
fn formatar(valor, casas) {
  let c = if casas == nil { 2 } else { casas }
  // ...
}
\end{lstlisting}

\section{Escopo de funções}

Funções declaradas com \lstinline{fn} no nível de módulo ficam disponíveis em todo
o arquivo. Funções declaradas dentro de blocos são locais ao bloco:

\begin{lstlisting}
{
  fn local(x) { x + 1 }
  display local(5)    // 6
}
display local(5)    // error: undefined function 'local'
\end{lstlisting}

\section{\texttt{const} dentro de funções}

\lstinline{const} pode ser usado dentro de funções para fixar valores que não devem
mudar durante a execução:

\begin{lstlisting}
fn calcular_imc(peso, altura) {
  const FORMULA = "peso / altura^2"
  peso / altura^2
}
\end{lstlisting}
